% TREŚĆ ROZDZIAŁU: usuń teksty instruktażowe i wpisz własne akapity.
% Nie zmieniaj \chapter ani \label, jeśli inne rozdziały odwołują się do tej
% etykiety. Nowe podrozdziały dodawaj przez \section{Tytuł}.
\chapter{Wprowadzenie}
\label{chap:wprowadzenie}

% Usuń wskazówki i zastąp je treścią projektu.
\begin{quote}\itshape
Krótko przedstaw kontekst, potrzebę realizacji projektu i problem, który ma zostać rozwiązany. Wyjaśnij znaczenie tematu dla odbiorcy. Nie opisuj tu jeszcze szczegółowo wymagań ani wyników.
\end{quote}

\section{Kontekst i motywacja}
Opisz dziedzinę, użytkownika lub proces, którego dotyczy projekt. Twierdzenia techniczne i dane pochodzące z literatury poprzyj cytowaniami, np. \cite{lamport1994}.

\section{Problem projektowy}
Zdefiniuj rzeczywistą trudność, a nie tylko czynność do wykonania. Wskaż stan początkowy i pożądany rezultat.

\section{Struktura dokumentu}
Jednym akapitem wyjaśnij rolę kolejnych rozdziałów raportu.
